\documentclass[12pt]{article}

\usepackage{amssymb,amsmath,amsthm}
\usepackage{ytableau}

% gaya teorema (saya menyukai semua lingkungan memakai pencacah yang sama); adaptasi berlisensi CC BY 4.0
\theoremstyle{plain}
\newtheorem{theorem}{Teorema}[section]
\newtheorem{lemma}[theorem]{Lema}
\newtheorem{proposition}[theorem]{Proposisi}
\newtheorem{corollary}[theorem]{Korolari}
\newtheorem{conjecture}[theorem]{Konjektur}
\newtheorem{exercise}[theorem]{Latihan}
\theoremstyle{definition}
\newtheorem{definition}[theorem]{Definisi}
\newtheorem{example}[theorem]{Contoh}
\theoremstyle{remark}
\newtheorem{remark}[theorem]{Catatan}

% beberapa pengaturan penomoran
\numberwithin{equation}{section}

\usepackage{fancyhdr}
\usepackage{lastpage}
\setlength{\headheight}{15.2pt}
\renewcommand{\footrulewidth}{0.4pt}% nilai bakunya 0pt
\setlength{\footskip}{30pt}
\pagestyle{fancy}

\makeatletter
\let\ps@plain\ps@fancy 
\makeatother

\lhead{MATH 742}
\chead{$S_n$ dan $\operatorname{GL}_n$}
\rhead{Musim Semi 2023}
\lfoot{Revisi Terakhir: \today}
\cfoot{}
\rfoot{\thepage\ dari \pageref{LastPage} }

\begin{document}

%%%%%%%%%%%%%%%%%%
%%%%%%%%%%%%%%%%%%
\title{Grup Simetrik dan Grup Linear Umum}
\author{Alexander Duncan}

\maketitle

Di sini kita membahas teori fungsi simetris, khususnya dengan
tujuan mendeskripsikan representasi grup simetrik dan grup
linear umum.
Representasi tak tereduksi dari grup simetrik
$S_n$ adalah \emph{modul Specht} $V_\lambda$, yang diparameterkan
oleh partisi $\lambda$ berbobot $n$.
Representasi polinomial tak tereduksi dari grup linear umum
$\operatorname{GL}(\mathbb{V})$ tepat merupakan (citra dari)
\emph{funktor Schur} $\mathbb{S}_\lambda(V)$, dengan $\lambda$
partisi yang panjangnya $\ell(\lambda) \le \dim(V)$.
Keduanya paling baik dipahami melalui bijeksi dengan \emph{fungsi
Schur} $\{ s_\lambda \}$, yang membentuk basis ortonormal bagi
gelanggang fungsi simetris. 

Catatan ini tidak berdiri sendiri. Banyak bukti hanya akan dibuat sketsanya atau
diserahkan sebagai rujukan. Alasannya bukan karena bukti-bukti itu sulit
(keindahan bidang ini justru terletak pada bukti yang sering sangat cerdik!), melainkan karena
teorinya terlalu kaya untuk dijelajahi secara memadai hanya dalam beberapa minggu.
Sebagian besar materi ini diambil dari
\cite[\S{4,6,A}]{FultonHarris},
\cite[\S{5.12--5.19}]{Etingof},
\cite[\S{I}]{Macdonald}, dan
\cite[\S{7}]{Stanley2}.

%%%%%%%%%%%%%%%%%%
%%%%%%%%%%%%%%%%%%

\section{Partisi}

Ingat bahwa suatu \emph{partisi} $\lambda$ adalah barisan bilangan bulat tak negatif
yang menurun lemah,
$\lambda_1 \ge \lambda_2 \ge \cdots$, dan akhirnya bernilai $0$.
Beberapa istilah baku:
\begin{itemize}
\item Nilai tak nol $\lambda_i$ disebut \emph{bagian} dari $\lambda$.
\item Banyaknya bagian, $\ell(\lambda)$, disebut \emph{panjang} dari $\lambda$.
\item Jumlah $|\lambda|=\sum_{i \ge 0} \lambda_i$ disebut \emph{bobot} dari $\lambda$.
\item ``Partisi dari $n$'' berarti partisi berbobot $n$.
\item $\lambda \vdash n$ berarti bahwa $\lambda$ adalah partisi dari $n$.
\item Banyaknya bagian $m_i(\lambda)$ yang sama dengan $i$ disebut
\emph{multiplisitas} $i$ dalam $\lambda$.
\item Partisi dapat ditulis sebagai $\lambda_1 + \lambda_2 + \cdots + \lambda_r$.
\item Kita dapat memakai notasi ringkas $\lambda = (1^{m_1} 2^{m_2} \cdots r^{m_r})$.
\item $\lambda+\mu$ adalah partisi $(\lambda_1+\mu_1, \ldots)$.
\item $\lambda \subseteq \mu$ berarti $\lambda_i \le \mu_i$ untuk semua $i \ge 1$.
\item Urutan dominasi $\lambda \le \mu$ berarti
$\sum_{j=1}^i\lambda_j \le \sum_{j=1}^i\mu_j$ untuk semua $i \ge 1$.
\end{itemize}

Partisi sering digambar sebagai \emph{diagram Young}. Diagram ini berupa
susunan kotak kosong dengan setiap baris berisi $\lambda_i$ kotak.
Kita memakai \emph{notasi Inggris}, dengan $\lambda_1$ sebagai baris teratas.
Konvensi di berbagai bidang matematika tidak selalu sama, jadi
bacalah setiap sumber dengan cermat.

Untuk suatu partisi $\lambda$, \emph{partisi konjugat}
adalah partisi $\lambda^\dag$ yang diperoleh dengan mencerminkan diagram Young terhadap
garis diagonal yang menurun ke kanan. Lebih eksplisit,
$\lambda^\dag_i$ adalah banyaknya bagian yang memenuhi $\lambda_j \ge i$. 

\begin{example}
Partisi-partisi dari $4$ adalah sebagai berikut:
\begin{center}
\begin{tabular}{ccccc}
4 & 3+1 & 2+2 & 2+1+1 & 1+1+1+1 \\
\ydiagram{4} & \ydiagram{3,1} & \ydiagram{2,2} & \ydiagram{2,1,1} &
\ydiagram{1,1,1,1}
\end{tabular}
\end{center}
\end{example}

\begin{example}
Misalkan $\lambda = (4,4,2,2,2,1)$.
Kita juga dapat menulis $\lambda$ sebagai $4+4+2+2+2+1$ atau $1^12^34^2$.
Bagian-bagiannya adalah $\lambda_1=4$, $\lambda_2=4$,
$\lambda_3=2$,
$\lambda_4=2$,
$\lambda_5=2$, dan
$\lambda_6=1$.
Panjangnya $\ell(\lambda)=6$, bobotnya $|\lambda|=15$,
dan multiplisitasnya $m_1(\lambda)=1$, $m_2(\lambda)=3$,
$m_3(\lambda)=0$, serta $m_4(\lambda)=2$.
Diagram Young-nya adalah
\begin{center}
\ydiagram{4,4,2,2,2,1}.
\end{center}
Partisi konjugat $\lambda^\dag$ adalah $(6,5,2,2)$.
\end{example}

Ingat bahwa setiap permutasi $\sigma \in S_n$ dapat ditulis sebagai
hasil kali siklus-siklus saling lepas, secara tunggal hingga pengurutan ulang
siklus-siklus tersebut. Dengan turut menghitung siklus panjang $1$, orde
siklus-siklus penyusunnya memberikan partisi $\lambda \vdash n$.
Sebagai contoh, $(1\ 4\ 3)(2\ 8)(6\ 7) \in S_9$ mempunyai tipe siklus
$3+2+2+1+1$.
Hasil berikut merupakan hasil baku dalam sebagian besar buku teori grup:

\begin{proposition}
Dua permutasi dalam $S_n$ saling konjugat jika dan hanya jika keduanya mempunyai
tipe siklus yang sama.
Khususnya, kelas-kelas konjugasi dari $S_n$ berkorespondensi bijektif secara kanonik
dengan partisi-partisi dari $n$.
\end{proposition}

Untuk setiap partisi $\lambda$, kita mendefinisikan bilangan bulat
\[
\epsilon_\lambda = (-1)^{|\lambda|-\ell(\lambda)} .
\]
dan
\[
z_\lambda = \prod_{i \ge 1} i^{m_i} m_i!
\]
dengan $m_i=m_i(\lambda)$ menyatakan multiplisitas,

\begin{exercise}
Buktikan bahwa jika $\sigma \in S_n$ mempunyai tipe siklus $\lambda$,
maka $\operatorname{sgn}(\sigma)=\epsilon_\lambda$.
\end{exercise}

\begin{exercise}
Buktikan bahwa jika $\sigma \in S_n$ mempunyai tipe siklus $\lambda$,
maka pemusat $Z_{S_n}(\sigma)$ berorde $z_\lambda$.
Secara ekuivalen, banyaknya elemen $S_n$ bertipe siklus $\lambda$
adalah $n!/z_\lambda$.
\end{exercise}

\section{Polinomial Simetris}

Terdapat aksi kiri alami dari grup simetrik $S_n$ pada gelanggang
polinomial $R=\mathbb{Z}[x_1,\ldots,x_n]$ dengan mempermutasikan peubah.
Lebih tepatnya, terdapat automorfisme gelanggang tunggal dari $R$ yang didefinisikan
oleh $x_i \mapsto x_{\sigma(i)}$ untuk setiap peubah $x_i$ dan setiap
permutasi $\sigma \in S_n$.
Sebagai rumusan lain, jika $\sigma \in S_n$, maka
\[
(\sigma \cdot f)(a_1,\ldots,a_n) =
f\left(a_{\sigma(1)}, \ldots, a_{\sigma}(n)\right)
\]
untuk $f \in R$ dan $a_1,\ldots,a_n \in \mathbb{Z}$.

\begin{definition}
Suatu polinomial $f \in \mathbb{Z}[x_1,\ldots, x_n]$ disebut \emph{simetris}
jika $\sigma \cdot f = f$ untuk semua $\sigma \in \Sigma$.
Kita menuliskan $\mathsf{SF}_n$ untuk himpunan polinomial simetris
$\mathbb{Z}[x_1,\ldots,x_n]^{S_n}$ dalam $n$ peubah.
\end{definition}

Himpunan $\mathsf{SF}_n$ membentuk gelanggang bergradasi
\[
\mathsf{SF}_n = \bigoplus_{d \ge 0} \mathsf{SF}_n^d
\]
dengan $\mathsf{SF}_n^d = \mathbb{Z}[x_1,\ldots,x_n]_{(d)}^{S_n}$
subgrup polinomial simetris homogen berderajat $d$.

Untuk suatu tupel $\alpha = (\alpha_1,\ldots,\alpha_n) \in \mathbb{N}^n$,
kita memakai notasi ringkas
\[
x^\alpha = x_1^{\alpha_1} \cdots x_n^{\alpha_n}
\]
untuk menyatakan monomial dalam $\mathbb{Z}[x_1,\ldots,x_n]$.

Untuk suatu partisi $\lambda \vdash d$ yang panjangnya $\ell(\lambda) \le n$,
\emph{polinomial simetris monomial yang terkait dengan $\lambda$}
diberikan oleh
\[
m_\lambda := \sum_{\alpha} x^\alpha
\]
dengan $(\alpha_1,\ldots,\alpha_n)$ merentang semua permutasi
\emph{berbeda} dari $(\lambda_1,\ldots,\lambda_n)$.

\begin{exercise}
Buktikan bahwa $\displaystyle z_\lambda m_\lambda = \sum_{\sigma \in S_n}
x^{\sigma(\lambda)}$ untuk semua $\lambda$ dengan panjang $\ell(\lambda) \le n$.
\end{exercise}

\begin{example}
Jika $n=3$, maka
\begin{align*}
m_\emptyset &= 1\\
m_1 &= x_1 + x_2 + x_3\\
m_2 &= x_1^2 + x_2^2 + x_3^2\\
m_{11} &= x_1x_2 + x_1x_3 + x_2x_3\\
m_{111} &= x_1x_2x_3 \\
m_{14} &= x_1x_2^4 + x_1^4x_2 + x_1x_3^4 + x_1^4x_3 + x_2x_3^4 +
x_2^4x_3.
\end{align*}
\end{example}

Yang penting, kita mempunyai hasil berikut:

\begin{proposition}
Polinomial simetris monomial membentuk basis bagi $\mathsf{SF}_n$.
Secara khusus,
\[
\mathsf{SF}_n^d = \operatorname{span}_{\mathbb{Z}}
\{ m_\lambda \mid \lambda \vdash d, \ell(\lambda) \le n \}.
\]
\end{proposition}

\begin{proof}
Partisi dengan panjang paling besar $n$ membentuk sistem wakil yang berbeda
bagi orbit-orbit $S_n$ dari $\mathbb{N}^n$.
Monomial-monomial dalam $m_\lambda$ tepat merupakan monomial pada orbit
$\lambda$.
Koefisien setiap monomial
dalam $m_\lambda$ adalah $0$ atau $1$, dan setiap monomial yang mungkin
muncul dalam tepat satu $m_\lambda$.
Jika $f$ polinomial simetris, maka koefisien
monomial $x^\alpha$ harus sama dengan koefisien $x^{\sigma(\alpha)}$
untuk setiap $\sigma \in S_n$.  
\end{proof}

Ada beberapa keluarga penting lain dari polinomial simetris
yang sekarang kita definisikan.

\begin{definition}
Untuk bilangan bulat positif $d$, kita mendefinisikan
\emph{polinomial simetris elementer $e_d$ berderajat $d$} sebagai
\[
e_d = m_{1^d} = \sum_{1 \le i_1 < i_2 < \cdots < i_d \le n} x_{i_1}x_{i_2}\cdots
x_{i_d},
\]
\emph{polinomial simetris homogen lengkap $h_d$ berderajat $d$} sebagai
\[
h_d = \sum_{\lambda \vdash d} m_\lambda = \sum_{1 \le i_1 \le i_2 \le \cdots \le i_d \le n} x_{i_1}x_{i_2}\cdots
x_{i_d},
\]
dan \emph{polinomial simetris jumlah pangkat $p_d$ berderajat $d$} sebagai
\[
p_d = m_d = \sum_{1 \le i \le n} x_i^d.
\]
Menurut konvensi, $e_0=h_0=p_0=1$.
\end{definition}

\begin{remark}
Konvensi untuk $e_0$ dan $h_0$ tidak menimbulkan perdebatan,
tetapi banyak rujukan membiarkan $p_0$ tak terdefinisi.
Memang ada alasan baik untuk mendefinisikan $p_0=n$, tetapi definisi ini tidak
meluas ke gelanggang fungsi simetris yang dibahas di bawah. 
\end{remark}

\begin{example}
Jika $n=3$, maka
\begin{align*}
e_1 = h_1 = p_1 &= x_1 + x_2 + x_3\\
e_2 &= x_1x_2 + x_1x_3 + x_2x_3\\
e_3 &= x_1x_2x_3 \\
e_4 &= 0 \\
h_2 &= x_1^2 + x_2^2 + x_3^2 + x_1x_2 + x_1x_3 + x_2x_3\\
h_3 &= x_1^3 + \cdots + x_1^2x_2 + \cdots + x_1x_2x_3\\
p_2 &= x_1^2 + x_2^2 + x_3^2\\
p_3 &= x_1^3 + x_2^3 + x_3^3
\end{align*}
\end{example}

Patut disebutkan sejak awal (meskipun buktinya kita tunda)
beberapa hasil mendasar. Pertama, kita mempunyai \emph{identitas Newton}:
\[
de_d = \sum_{i=1}^d (-1)^{i-1} p_i e_{d-i} ,
\]
serta relasi mendasar
\[
0 = \sum_{i=0}^d (-1)^i e_i h_{d-i},
\]
yang berlaku untuk semua bilangan bulat positif $d$.
Relasi ini memungkinkan kita menghitung $p_i$ dan $h_i$ secara rekursif dalam
$e_i$ (dan sebaliknya).

Kita juga mempunyai \emph{Teorema Dasar Polinomial
Simetris}, yang menyatakan bahwa $e_1,\ldots,e_n$ adalah pembangkit bebas secara aljabar
dari gelanggang $\mathsf{SF}_n$.
Bahkan, dengan menggabungkan hasil ini dan relasi di atas, kita memperoleh
\begin{align*}
\mathbb{Z}[x_1,\ldots,x_n]^{S_n} &= \mathbb{Z}[e_1,\ldots,e_n]\\
\mathbb{Z}[x_1,\ldots,x_n]^{S_n} &= \mathbb{Z}[h_1,\ldots,h_n]\\
\mathbb{Q}[x_1,\ldots,x_n]^{S_n} &= \mathbb{Q}[p_1,\ldots,p_n]
\end{align*}
dengan koefisien rasional diperlukan untuk kesamaan terakhir.
Bukti kesamaan-kesamaan ini tidak terlalu sulit, tetapi kita menundanya
karena semuanya merupakan akibat dari fakta yang lebih umum.

\subsection{Polinomial Berselang-seling dan Polinomial Schur}

Misalkan $\epsilon : S_n \to \{\pm 1\}$ adalah homomorfisme tanda
dan $A_n = \ker(\epsilon)$ adalah grup berselang-seling pada $n$ huruf.

\begin{definition}
Polinomial $f \in \mathbb{Z}[x_1,\ldots,x_n]$ disebut \emph{berselang-seling} jika
$\sigma \cdot f = \epsilon(\sigma)f$ untuk semua $\sigma \in S_n$.
Untuk suatu elemen $\alpha \in \mathbb{N}^n$, definisikan \emph{alternan
dari $\alpha$} sebagai
\[
a_\alpha = \sum_{\sigma \in S_n} \epsilon(\sigma) x^{\sigma(\alpha)},
\]
yang merupakan polinomial berselang-seling.
\end{definition}

Perhatikan bahwa $a_\alpha=0$ jika dan hanya jika entri-entri $\alpha$
tidak semuanya berbeda. Kita juga melihat bahwa setiap polinomial berselang-seling
adalah kombinasi
linear dari $a_\lambda$ untuk partisi $\lambda$ yang semua bagiannya
$\lambda_i$ berbeda.
Perhatikan bahwa semua bagian $\lambda$ berbeda jika dan hanya jika
$\delta \subseteq \lambda$, dengan
$\delta=(n-1,n-2,\ldots,1,0)$.

Kita mempunyai hasil berikut:
\begin{lemma}
Himpunan
\[
\{ a_\lambda \mid \ell(\lambda) \le n, \delta \subseteq \lambda \}
\]
merupakan basis bagi grup polinomial berselang-seling.

\end{lemma}
\begin{example}
\emph{Polinomial Vandermonde} $\Delta \in \mathbb{Z}[x_1,\ldots,x_n]$
didefinisikan oleh
\[
\Delta = \prod_{1 \le i < j \le n} (x_i-x_j).
\]
Perhatikan bahwa $(i\ j) \cdot \Delta = -\Delta$ untuk semua $i \ne j$.
Kita menyimpulkan bahwa $\sigma(\Delta) = \operatorname{sgn}(\sigma)$
untuk semua $\sigma \in S_n$.
Jadi $\Delta$ berselang-seling.
Perhatikan bahwa $\Delta=a_\delta$.

\end{example}
\begin{lemma}
Setiap polinomial berselang-seling $f$ berbentuk $g \Delta$
dengan $g$ polinomial simetris. 

\end{lemma}
\begin{proof}
Cukup ditunjukkan bahwa $(x_i-x_j)$ membagi $f$ untuk semua $i \ne j$.
Memang, andaikan $c_\alpha x^\alpha$ adalah monomial dalam $f$
dan tulis $x^\alpha = x_i^u x_j^v x^\beta$ dengan $x_i$ dan $x_j$
tidak membagi $x^\beta$.
Karena $f$ berselang-seling, $-c_\alpha x_i^v x_j^u x^\beta$
juga harus muncul dalam $f$.
Jadi $c_\alpha (x_i^u x_j^v - x_i^v x_j^u) x^\beta$ muncul dalam $f$.
Masing-masing suku ini habis dibagi oleh $(x_i-x_j)$ sebagaimana diinginkan.

\end{proof}

Definisi berikut kini bermakna:
\begin{definition}
Untuk partisi $\lambda$ dengan panjang $\le n$, \emph{polinomial Schur}
adalah polinomial simetris
$s_\lambda = a_{\lambda+\delta}/a_\delta$.

\end{definition}
Alternan $\{ a_{\lambda+\delta} \}$ membentuk basis bagi
polinomial berselang-seling, sebagaimana polinomial simetris monomial
$\{ m_\lambda \}$ membentuk basis bagi polinomial
simetris.
Karena $a_{\lambda+\delta}=s_\lambda a_\delta$, polinomial Schur dapat dipandang
sebagai ``polinomial monomial berselang-seling'', tetapi faktor

berselang-selingnya telah dibagi sehingga hasilnya simetris.

Khususnya, kita mempunyai hasil berikut:
\begin{proposition}
Himpunan polinomial Schur
\[
\{ s_\lambda \mid \lambda \vdash d, \ell(\lambda) \le n \}
\]
merupakan basis bagi $\mathsf{SF}^d_n$.

\end{proposition}
\emph{Bilangan Kostka} $K_{\lambda \mu}$ adalah entri-entri matriks
perubahan basis antara basis monomial dan basis Schur.
Secara khusus, bilangan-bilangan ini adalah bilangan bulat tak negatif sedemikian sehingga
\[
s_\lambda = \sum_{\substack{\mu\, \vdash |\lambda|\\\ell(\mu) \le n}} K_{\lambda \mu} m_\mu
\]
untuk semua partisi $\lambda,\mu$. Bilangan Kostka tidak bergantung pada banyaknya

peubah $n$ dalam gelanggang polinomial yang memuatnya
(hal ini tidak kita buktikan di sini).
\begin{example}
Untuk $n \ge 4$, kita mempunyai matriks perubahan basis berikut:
\[
\begin{pmatrix}
s_{1111}\\ s_{211} \\ s_{22}\\ s_{31}\\ s_4
\end{pmatrix}
=
\begin{pmatrix}
1 & 0 & 0 & 0 & 0\\
3 & 1 & 0 & 0 & 0\\
2 & 1 & 1 & 0 & 0\\
3 & 2 & 1 & 1 & 0\\
1 & 1 & 1 & 1 & 1
\end{pmatrix} 
\begin{pmatrix}
m_{1111}\\ m_{211} \\ m_{22}\\ m_{31}\\ m_4
\end{pmatrix}
\]

Jadi, misalnya, $K_{(31)(1111)}=3$.
\end{example}
\begin{exercise}

Buktikan bahwa $s_{(d)} = h_d$ dan $s_{(1^d)} = e_d$.

\end{exercise}
Kita tidak akan memerlukan hasil berikut nanti, tetapi hasil ini menarik secara umum:
\begin{theorem}[Teorema Dasar Polinomial Berselang-seling]
Gelanggang $\mathbb{Z}[x_1,\ldots,x_n]^{A_n}$ adalah modul-$\mathsf{SF}_n$
bebas dengan basis $\{1, \Delta\}$.
Dengan kata lain, setiap polinomial invarian-$A_n$ $f$ dapat ditulis secara tunggal

sebagai $f=p+\Delta q$, dengan $p,q$ simetris.
\end{theorem}
\begin{remark}
Jika $V$ representasi linear kompleks berdimensi $n$ dari $G$,
maka terdapat aksi alami $G$ pada gelanggang polinomial
$S=\mathbb{C}[x_1,\ldots,x_n]$ dengan memandangnya sebagai aljabar simetris
$\mathcal{S}(V^\vee)$.
Sebuah teorema terkenal Hochster--Roberts membuktikan bahwa gelanggang invarian
$S^G=\mathbb{C}[x_1,\ldots,x_n]^G$ adalah gelanggang Cohen--Macaulay:
terdapat subgelanggang polinomial $R=\mathbb{C}[f_1,\ldots,f_n]$
sedemikian sehingga $S^G$ merupakan modul-$R$ bebas.
``Teorema Dasar Polinomial Berselang-seling'' dapat dipandang sebagai
kasus yang sangat khusus dari fakta ini (dengan subgelanggang dan basis

yang mempunyai interpretasi tertentu).

\end{remark}
\subsection{Pangkat Simetris dan Eksterior}

Kini kita melihat beberapa kaitan antara fungsi simetris dan
teori representasi:
\begin{proposition}
Andaikan $V$ ruang vektor berdimensi $n$ dan
$\varphi \in \operatorname{End}(V)$ mempunyai nilai-nilai eigen
$\alpha_1,\ldots,\alpha_n$.
Maka aksi alami $\varphi$ pada pangkat eksterior $\Lambda^d V$
mempunyai teras
\[
\operatorname{tr}\left( \Lambda^d \varphi \right)
= e_d(\alpha_1,\ldots,\alpha_n)
\]
dan aksi alami $\varphi$ pada pangkat simetris
$\mathcal{S}^d V$ mempunyai teras
\[
\operatorname{tr}\left( \mathcal{S}^d \varphi \right)
= h_d(\alpha_1,\ldots,\alpha_n).
\]
\end{proposition}

\begin{proof}
Cukup bekerja atas medan tertutup secara aljabar $k$.

Kita membuktikan pernyataan untuk $\Lambda^d V$ karena argumen untuk
$\mathcal{S}^d$ sangat serupa. 
Pertama, andaikan $\varphi$ dapat didiagonalkan, dengan vektor-vektor eigen
$v_1,\ldots, v_n$ yang bersesuaian dengan nilai-nilai eigen
$\alpha_1,\ldots,\alpha_n$.
Suatu basis eigen bagi $\Lambda^d V$ diberikan oleh
\[
\{ v_{i_1}\wedge v_{i_2}\wedge \cdots \wedge v_{i_d} \mid
1 \le i_1 < i_2 < \cdots < i_d \le n \} \ .
\]
Jadi nilai-nilai eigen yang bersesuaian dari $\Lambda^d \varphi$ adalah
\[
\{ \alpha_{i_1} \alpha_{i_2} \cdots \alpha_{i_d} \mid
1 \le i_1 < i_2 < \cdots < i_d \le n \} \ .
\]

Teras $\Lambda^d \phi$ hanyalah jumlah nilai-nilai eigen tersebut, yakni
$e_d(\alpha_1,\ldots,\alpha_n)$ sebagaimana diinginkan.
Sekarang tinjau kasus ketika $\varphi$ tidak dapat didiagonalkan.
Dalam kasus ini $\varphi$ mempunyai bentuk kanonik Jordan. Jadi
terdapat basis $v_1,\ldots, v_n$ bagi $V$ sehingga $\varphi$
direpresentasikan oleh $D+N$, dengan $D$ matriks diagonal dan

$N$ matriks segitiga bawah.
Perhatikan bahwa $N(v_i)$ adalah kombinasi linear dari $v_{i+1},\ldots, v_n$.
Definisikan urutan total pada $\mathbb{N}^n$, dengan
$(a_1,\ldots,a_n) < (b_1,\ldots,b_n)$ jika $a_i < b_i$ untuk indeks terkecil $i$
yang memenuhi $a_i \ne b_i$.
Kini perhatikan bahwa
\[
(D+N)(v_{i_1}\wedge \cdots \wedge v_{i_d})
= D(v_{i_1}\wedge \cdots \wedge v_{i_d})
+ \textrm{suku-suku lebih tinggi} 
\]
dengan ``suku-suku lebih tinggi'' merupakan kelipatan skalar dari
$v_{j_1}\wedge \cdots \wedge v_{j_d}$
dan $(j_1,\ldots,j_d) > (i_1,\ldots,i_d)$.
Jadi, bila diurutkan secara tepat, basis kita bagi $\Lambda^d V$
juga merepresentasikan $\Lambda^d V$ sebagai matriks segitiga bawah.
Teras hanya bergantung pada entri diagonal, sehingga hasilnya mengikuti

dari kasus yang dapat didiagonalkan.
\end{proof}
Cara lain memahami proposisi sebelumnya adalah bahwa
$\Lambda^d V$ dan $\mathcal{S}^d V$ merupakan representasi dari
$\operatorname{GL}(V)$ dengan karakter masing-masing $e_d$ dan
$h_d$.
Kita akan melihat bahwa terdapat suatu kelas polinomial simetris
yang disebut \emph{polinomial Schur}, yang tepat merupakan karakter


dari representasi \emph{polinomial} tak tereduksi

$\operatorname{GL}(V)$.
\section{Gelanggang Fungsi Simetris}
Banyak relasi antara berbagai polinomial simetris khusus pada dasarnya
tidak bergantung pada banyaknya peubah $n$, yakni $x_1,\ldots,x_n$,
dalam gelanggang polinomial yang memuatnya. \emph{Gelanggang fungsi simetris}
adalah objek baku dalam kombinatorika aljabar yang memudahkan

pembahasan ini. Ada beberapa konstruksi yang ekuivalen; kita mengikuti
konstruksi dalam \cite{Stanley2}, yang cukup konkret.
Misalkan $\mathbb{Z}[[\{x_n\}_{n \in \mathbb{N}_{>0}}]]$ adalah gelanggang deret
pangkat formal dalam terhitung banyak peubah. Ingat bahwa ini berarti
setiap monomial hanya memuat hingga banyak $x_i$, tetapi setiap elemen boleh berupa
kombinasi linear tak hingga dari monomial, dan mungkin tidak ada batas atas
bagi subskrip $i$ yang muncul di antara $x_i$ dalam setiap
monomial.
Untuk bilangan bulat tak negatif $d$, subgrup
$\mathbb{Z}[[\{x_n\}_{n \in \mathbb{N}_{>0}}]]_{(d)}$ terdiri atas

elemen-elemen yang semua monomialnya tepat berderajat $d$ (dengan kata lain,
tepat $d$ peubah $x_i$ dalam setiap monomial, dihitung beserta multiplisitasnya).
Misalkan $S_{\mathbb{N}_{>0}}$ adalah grup simetrik pada $\mathbb{N}_{>0}$; dengan
kata lain, grup semua bijeksi $\mathbb{N}_{>0} \to \mathbb{N}_{>0}$.
Terdapat aksi alami $S_{\mathbb{N}_{>0}}$ pada
$\mathbb{Z}[[\{x_n\}_{n \in \mathbb{N}_{>0}}]]$ dengan mempermutasikan peubah, dan aksi ini
mempertahankan derajat.
Suatu \emph{fungsi simetris homogen berderajat $d$} adalah elemen
$f \in \mathbb{Z}[[\{x_n\}_{n \in \mathbb{N}_{>0}}]]_{(d)}^{S_{\mathbb{N}}}$.

Misalkan $\mathsf{SF}^d$ adalah subgrup
fungsi simetris homogen berderajat $d$.
Kita mendefinisikan \emph{gelanggang fungsi simetris} sebagai jumlah langsung (internal)
\[
\mathsf{SF} = \bigoplus_{d \ge 0} \mathsf{SF}^d,
\]

yang merupakan subgelanggang bergradasi dari $\mathbb{Z}[[\{x_n\}_{n \in \mathbb{N}_{>0}}]]$.
\begin{remark}
Perhatikan bahwa ``fungsi simetris'' merupakan istilah baku,
tetapi bukan nama yang ideal karena tidak benar-benar jelas
apa yang menjadi argumen \emph{fungsi} tersebut.
Selain itu, istilah fungsi simetris sering digunakan dalam kasus
berhingga banyak peubah untuk membahas fungsi ``biasa'' seperti $e^{x+y}$,

yang invarian terhadap simetri seperti $x \leftrightarrow y$.
\end{remark}
Terdapat homomorfisme gelanggang bergradasi surjektif kanonik
$\rho_n : \mathsf{SF} \to \mathsf{SF}_n$, yang didefinisikan oleh
\[
\rho_n(f)(x_1,\ldots,x_n) = f(x_1,\ldots,x_n,0,0,\cdots) .
\]

Perhatikan bahwa konvergensi bukan masalah, karena semua kecuali hingga banyak monomial
akan bernilai nol.
\begin{remark}
Sebagai rumusan ekuivalen, gelanggang fungsi simetris dapat didefinisikan sebagai limit
invers
\[
\mathsf{SF} = \lim_{\substack{\longleftarrow\\n}} \mathsf{SF}^n
\]
dalam kategori gelanggang bergradasi (lihat \cite{Macdonald}).
(Peringatan: ini \emph{bukan} limit invers dalam kategori gelanggang biasa.)
Pemetaan $\rho_n$ di atas tepat merupakan proyeksi kanonik yang diperoleh dari

konstruksi ini.
\end{remark}
Kita mendefinisikan \emph{fungsi simetris monomial yang terkait dengan $\lambda$}
sebagai
\[
m_\lambda := \sum_{\alpha} x^\alpha
\]
dengan $(\alpha_1,\ldots) \in \mathbb{N}^{\mathbb{N}_{>0}}$ merentang semua permutasi

\emph{berbeda} dari $(\lambda_1,\ldots)$.
Perhatikan bahwa himpunan $\{ m_\lambda\}$ merupakan basis bagi $\mathsf{SF}$.
Sekarang kita dapat mendefinisikan \emph{fungsi simetris elementer}
\[
e_d = m_{1^d},
\]
\emph{fungsi simetris homogen lengkap}
\[
h_d = \sum_{\lambda \vdash d} m_\lambda ,
\]
dan \emph{fungsi simetris jumlah pangkat}
\[
p_d = m_{d}.
\]

Pemakaian ulang notasi $m_\lambda$, $e_d$, $h_d$, dan $p_d$ untuk menyatakan
objek berbeda dalam setiap $\mathsf{SF}^n$ pada umumnya tidak menimbulkan masalah,

mengingat semuanya tepat merupakan citra objek yang
bersesuaian dalam $\mathsf{SF}$ di bawah pemetaan ke $\rho_n$.
Kita juga dapat mendefinisikan \emph{fungsi simetris Schur}
$s_\lambda$ melalui 
\[
s_\lambda = \sum_{\mu} K_{\lambda \mu} m_\mu
\]

dengan memanfaatkan fakta bahwa bilangan Kostka

tidak bergantung pada banyaknya peubah dalam gelanggang polinomial yang diberikan.
\subsection{Relasi dan Identitas}
\begin{proposition} \label{prop:fund_rel}

$\displaystyle\sum_{i=0}^d (-1)^i e_i h_{d-i}=0$ untuk setiap $d \ge 1$.
\end{proposition}
\begin{proof}
Tinjau fungsi pembangkit
\[ E(t) = \sum_{r \ge 0} e_r t^r
\textrm{ dan }
H(t) = \sum_{r \ge 0} h_r t^r \]
sebagai elemen dalam $\mathsf{SF}[[t]]$.
Kita melihat bahwa
\begin{align*}
& E(t)\\
 =& 1 + (x_1+x_2+\cdots)t
+ (x_1x_2+x_1x_3+\cdots)t^2 \\
=&\prod_{i \ge 1} (1+x_i t)
\end{align*}
dalam gelanggang yang lebih besar $\mathbb{Z}[[x_1,\ldots]][[t]]$.
Serupa dengan itu, dengan mengingat rumus deret geometri, kita mempunyai
\[
H(t) = \prod_{i \ge 1} (1-x_i t)^{-1} .
\]
Kesamaan $H(t)E(-t)=1$ diperoleh langsung.
Identitas yang diinginkan hanyalah koefisien $t^d$

dalam kesamaan tersebut.
\end{proof}
\begin{proposition}[Identitas Newton] Untuk semua $d \ge 1$,
\begin{align*}
dh_d&=\sum_{i=1}^d p_i h_{d-i} \textrm{ dan} \\ 
de_d&=\sum_{i=1}^d (-1)^{i-1} p_i e_{d-i}.
\end{align*}

\end{proposition}
\begin{proof}
Kita menghitung
\begin{align*}
\sum_{r \ge 1} p_r t^{r-1}
&= \sum_{i \ge 1} \sum_{d \ge 0} x_i^d t^{d-1}
= \sum_{i \ge 1} \frac{x_i}{1-x_it}\\
&= \sum_{i \ge 1} \frac{d}{dt}\log\left(
\left(1-x_i\right)^{-1}\right)
= \frac{H'(t)}{H(t)} = \frac{E'(-t)}{E(-t)}
\end{align*}
kemudian identitasnya dibaca dari
koefisien $t^d$ dalam $P(t)H(t)=H'(t)$

dan $P(t)E(-t)=E'(-t)$.
\end{proof}
Untuk $\lambda=(\lambda_1,\lambda_2,\ldots,\lambda_r)$ kita mendefinisikan
$e_\lambda = e_{\lambda_1}\cdot e_{\lambda_2} \cdots e_{\lambda_r}$,

$h_\lambda = h_{\lambda_1}\cdot h_{\lambda_2} \cdots h_{\lambda_r}$, dan
$p_\lambda = p_{\lambda_1}\cdot p_{\lambda_2} \cdots p_{\lambda_r}$.
Ingat \emph{urutan parsial dominasi} pada partisi, dengan $\lambda \ge \mu$
jika dan hanya jika
$\lambda_1 + \cdots + \lambda_i \ge \mu_1 + \cdots \mu_i$
untuk semua $i \ge 1$.

(Perhatikan bahwa ini hanya urutan parsial; sebagai contoh, $(31^3)$ dan $(2^3)$
tidak dapat dibandingkan.
\begin{proposition} \label{prop:01matrix}
Kita mempunyai
\[
e_\lambda = \sum_{\mu \vdash |\lambda|} M_{\lambda \mu} m_\mu
\]
dengan $M_{\lambda \mu}$ banyaknya matriks-$\{0,1\}$ yang jumlah barisnya
berturut-turut $\lambda_1,\lambda_2,\ldots$ dan jumlah kolomnya
berturut-turut $\mu_1,\mu_2,\ldots$.
Selain itu, $M_{\lambda\mu} \ne 0$ jika dan hanya jika $\mu \le \lambda^\dag$

dan $M_{\lambda\lambda^\dag}=1$.
\end{proposition}
\begin{proof}
Kita hanya membuat sketsa argumennya.
Kita meninjau suku-suku $e_\lambda$ dan $m_\mu$ dengan mengacu pada
matriks berikut:
\[
X = \begin{pmatrix}
x_1 & x_2 & x_3 & \cdots \\
x_1 & x_2 & x_3 & \cdots \\
x_1 & x_2 & x_3 & \cdots \\
\vdots & \vdots & \vdots & \ddots
\end{pmatrix}.
\]
Perhatikan bahwa setiap matriks-$\{0,1\}$ $Y$ menghasilkan monomial $x^\alpha$
dengan mengambil hasil kali entri-entri $X$ yang bersesuaian dengan
entri tak nol dari $Y$.
Monomial dalam $e_\lambda$ dikonstruksi secara tunggal dengan
mengambil hasil kali tepat $\lambda_1$ entri berbeda dari baris $1$,
tepat $\lambda_2$ dari baris $2$, dan seterusnya.
Monomial $x^\mu$ dikonstruksi dengan

mengambil hasil kali tepat $\mu_1$ entri berbeda dari kolom $1$,
tepat $\mu_2$ dari kolom $2$, dan seterusnya.
Syarat pada $M_{\lambda\mu}$ kini diperoleh dengan mencari

matriks-$\{0,1\}$ yang memenuhi kendala-kendala tertentu.
\end{proof}
Setelah memperhalus urutan parsial kita menjadi urutan total
pada partisi, $M=\left(M_{\lambda\mu}\right)$ dapat dipandang sebagai
matriks bilangan bulat tak hingga.

Perhatikan bahwa $M_{\lambda\mu^\dag}$ berbentuk segitiga dengan $1$ pada diagonal.
Jadi $M$ dapat dibalik dan kita memperoleh korolari penting berikut:
\begin{corollary}[Teorema Dasar Fungsi Simetris]
Terdapat kesamaan gelanggang
\[\mathsf{SF} = \mathbb{Z}[e_1,e_2,\ldots] \]

dengan $e_1,e_2,\ldots$ bebas secara aljabar.
\end{corollary}
\begin{exercise}
Carilah analog Proposisi~\ref{prop:01matrix} untuk $h_\lambda$
dengan memakai matriks-$\mathbb{N}$ alih-alih matriks-$\{0,1\}$.
Simpulkan bahwa 
$\mathsf{SF} = \mathbb{Z}[h_1,h_2,\ldots]$

dengan $h_1,h_2,\ldots$ bebas secara aljabar.
\end{exercise}
\begin{exercise}
Gunakan identitas Newton untuk menunjukkan bahwa 
$\mathsf{SF} \otimes_{\mathbb{Z}} \mathbb{Q} \cong \mathbb{Q}[p_1,p_2,\ldots]$

dengan $p_1,p_2,\ldots$ bebas secara aljabar.
\end{exercise}
\begin{example}
\[
\begin{pmatrix} e_{11}\\ e_2\\ e_{111}\\ e_{21}\\ e_3\\
e_{1111}\\ e_{211} \\ e_{22}\\ e_{31}\\ e_4
\end{pmatrix}
=
\begin{pmatrix}
2 & 1 &
 & & &
& & & & \\
1 & &
 & & &
& & & & \\
 & &
6 & 3 & 1 &
& & & & \\
 & &
3 & 1 & &
& & & & \\
 & &
1 & & &
& & & & \\
 & &
 & & &
24 & 12 & 6 & 4 & 1 \\
 & &
 & & &
12 & 5 & 2 & 1 & \\
 & &
 & & &
6 & 2 & 1 & & \\
 & &
 & & &
4 & 1 & & & \\
 & &
 & & &
1 & & & &
\end{pmatrix} 
\begin{pmatrix} m_{11}\\ m_2\\ m_{111}\\ m_{21}\\ m_3\\
m_{1111}\\ m_{211} \\ m_{22}\\ m_{31}\\ m_4
\end{pmatrix}
\]

\end{example}
\begin{example}
\[
\begin{pmatrix} h_{11}\\ h_2\\ h_{111}\\ h_{21}\\ h_3\\
h_{1111}\\ h_{211} \\ h_{22}\\ h_{31}\\ h_4
\end{pmatrix}
=
\begin{pmatrix}
2 & 1 &
 & & &
& & & & \\
1 & 1 &
 & & &
& & & & \\
 & &
6 & 3 & 1 &
& & & & \\
 & &
3 & 2 & 1 &
& & & & \\
 & &
1 & 1 & 1 &
& & & & \\
 & &
 & & &
24 & 12 & 6 & 4 & 1 \\
 & &
 & & &
12 & 7 & 4 & 3 & 1 \\
 & &
 & & &
6 & 4 & 3 & 2 & 1 \\
 & &
 & & &
4 & 3 & 2 & 2 & 1 \\
 & &
 & & &
1 & 1 & 1 & 1 & 1
\end{pmatrix} 
\begin{pmatrix} m_{11}\\ m_2\\ m_{111}\\ m_{21}\\ m_3\\
m_{1111}\\ m_{211} \\ m_{22}\\ m_{31}\\ m_4
\end{pmatrix}
\]

\end{example}
\begin{example}
\[
\begin{pmatrix} e_{11}\\ e_2\\ e_{111}\\ e_{21}\\ e_3\\
e_{1111}\\ e_{211} \\ e_{22}\\ e_{31}\\ e_4
\end{pmatrix}
=
\begin{pmatrix}
1 & &
 & & &
& & & & \\
\frac{1}{2} & -\frac{1}{2} &
 & & &
& & & & \\
 & &
1 & & &
& & & & \\
 & &
\frac{1}{2} & -\frac{1}{2} &  &
& & & & \\
 & &
\frac{1}{6} & -\frac{1}{2} & \frac{1}{3} &
& & & & \\
 & &
 & & &
1 & & & & \\
 & &
 & & &
\frac{1}{2} & -\frac{1}{2} & & & \\
 & &
 & & &
\frac{1}{4} & -\frac{1}{2} & \frac{1}{4} & & \\
 & &
 & & &
\frac{1}{6} & -\frac{1}{2} & & \frac{1}{3} & \\
 & &
 & & &
\frac{1}{24} & -\frac{1}{4} & \frac{1}{8} & \frac{1}{3} & -\frac{1}{4}
\end{pmatrix} 
\begin{pmatrix}
p_{11}\\ p_2\\ p_{111}\\ p_{21}\\ p_3\\
p_{1111}\\ p_{211} \\ p_{22}\\ p_{31}\\ p_4
\end{pmatrix}
\]


\end{example}
\subsection{Tableau Young}
Sekarang kita berikan interpretasi kombinatorial bagi fungsi Schur
dan bilangan Kostka.
Suatu \emph{tableau Young semistandar $T$ berbentuk $\lambda$}
adalah diagram Young dari partisi $\lambda$ yang kotak-kotaknya diisi
bilangan bulat positif yang meningkat lemah pada setiap baris
dan meningkat ketat pada setiap kolom.
\emph{Tipe} atau \emph{isi} $\alpha$ dari tableau Young semistandar
$T$ adalah barisan
$(\alpha_1,\alpha_2,\ldots)$ bilangan bulat tak negatif
dengan $\alpha_i$ banyaknya kotak yang berisi $i$.
Untuk tableau Young semistandar $T$ bertipe $\alpha$,
kita mempunyai monomial
\[
x^T = x_1^{\alpha_1}x_2^{\alpha_2} \cdots .
\]

\begin{example}
Berikut adalah tableau Young berbentuk $422$
dan bertipe $(2,3,0,1,2)$:
\begin{center}
\begin{ytableau}
1 & 1 & 2 & 5 \\
2 & 2 \\
4 & 5
\end{ytableau}
\end{center}

Monomial yang bersesuaian adalah $x^T=x_1^2x_2^3x_4x_5^2$.
\end{example}
Hasil berikut dipakai sebagai \emph{definisi} fungsi Schur dalam

\cite{Stanley2}; kesesuaiannya dengan definisi klasik
di atas diberikan oleh \cite[Theorem 7.15.2]{Stanley2}.
\begin{theorem}
Jika $\lambda$ suatu partisi, maka
\[
s_\lambda = \sum_{T} x^T
\]
dengan jumlah diambil atas semua tableau Young semistandar berbentuk
$\lambda$.
Khususnya, $K_{\lambda \mu}$ adalah banyaknya tableau Young semistandar

berbentuk $\lambda$ dan bertipe $\mu$.
\end{theorem}

Berbagai sifat bilangan Kostka lebih mudah diperiksa
dengan definisi ini:
\begin{exercise}
Untuk partisi $\lambda, \mu$, berlaku
$K_{\lambda \mu}=0$, kecuali jika $\mu \le \lambda$ dalam urutan dominasi.

Selain itu, $K_{\lambda \lambda}=1$.
\end{exercise}
Bahkan, hampir semua matriks transisi dari
$\mathsf{SF}$ dapat dipahami melalui bilangan Kostka
(lihat \cite[\S{I.6}]{Macdonald} untuk pembahasan lengkap).
Beberapa contoh penting:
\begin{align*}
s_\lambda &= \sum_{\mu} K_{\lambda \mu} m_\mu\\
h_\lambda &= \sum_{\mu} K_{\mu\lambda} s_\mu\\
e_\lambda &= \sum_{\mu} K_{\mu^\dag\lambda} s_\mu.
\end{align*}

Untuk partisi $\lambda \vdash n$, bilangan Kostka khusus
$f^\lambda = K_{\lambda(1)^n}$ sangat penting.
Tepatnya, $f^\lambda$ menghitung banyaknya \emph{tableau Young standar}
yang kotak-kotaknya masing-masing memuat bilangan dalam $\{1,\ldots,n\}$ tepat sekali
(dan tetap meningkat sepanjang baris dan kolom).
Untuk kotak ke-$(i,j)$ dari diagram Young, misalkan \emph{panjang kait}

$h(i,j)$ menghitung banyaknya kotak yang terletak tepat di bawah dan tepat
di kanan kotak $(i,j)$, beserta kotak itu sendiri.
\begin{theorem}[Rumus Panjang Kait] \label{thm:hook_length}
$f^\lambda = \frac{n!}{\prod h(i,j)}$

dengan hasil kali diambil atas semua kotak diagram Young dari $\lambda$.
\end{theorem}
\begin{example}
Tinjau partisi $\lambda = (4,3,2)$.
Kita mengisi setiap kotak diagram Young dengan panjang kaitnya:
\begin{center}
\begin{ytableau}
6 & 5 & 3 & 1 \\
4 & 3 & 1 \\
2 & 1
\end{ytableau}
\end{center}
Jadi,
\[
f^\lambda = \frac{9!}{6\cdot 5 \cdot 4 \cdot 3 \cdot 3 \cdot 2}
= 168
\]

dalam kasus ini.
\end{example}
Perubahan basis antara jumlah pangkat dan fungsi Schur

sangat menarik bagi teori representasi grup simetrik,
seperti akan segera kita lihat.
\begin{example} \label{ex:transSP}
\[
\begin{pmatrix}
p_{11}\\ p_2\\ p_{111}\\ p_{21}\\ p_3\\
p_{1111}\\ p_{211} \\ p_{22}\\ p_{31}\\ p_4
\end{pmatrix}
=
\begin{pmatrix}
1 & 1 &
 & & &
& & & & \\
-1 & 1 &
 & & &
& & & & \\
 & &
1 & 2 & 1 &
& & & & \\
 & &
-1 & 0 & 1 &
& & & & \\
 & &
1 & -1 & 1 &
& & & & \\
 & &
 & & &
1 & 3 & 2 & 3 & 1 \\
 & &
 & & &
-1 & -1 & 0 & 1 & 1 \\
 & &
 & & &
1 & -1 & 2 & -1 & 1 \\
 & &
 & & &
1 & 0 & -1 & 0 & 1 \\
 & &
 & & &
-1 & 1 & 0 & -1 & 1
\end{pmatrix} 
\begin{pmatrix} s_{11}\\ s_2\\ s_{111}\\ s_{21}\\ s_3\\
s_{1111}\\ s_{211} \\ s_{22}\\ s_{31}\\ s_4
\end{pmatrix}
\]


\end{example}
Kita menyebutkan beberapa konstruksi berguna lainnya (tanpa bukti)

yang semoga mendorong pembaca untuk
mendalami \cite{Macdonald} dan \cite{Stanley2}.
\emph{Hasil kali dalam Hall} adalah bentuk bilinear simetris
definit positif tunggal $(-,-)$ pada $\mathsf{SF}$ yang memenuhi
\begin{align*}
( s_\lambda, s_\mu ) &= \delta_{\lambda \mu},\\
( h_\lambda, m_\mu ) &= \delta_{\lambda \mu}, \textrm{ dan}\\
( p_\lambda, p_\mu ) &= z_\lambda \delta_{\lambda \mu}
\end{align*}

untuk semua partisi $\lambda, \mu$.
Salah satu akibat Teorema Dasar adalah adanya isomorfisme
gelanggang
\[
\omega : \mathsf{SF} \to \mathsf{SF}
\]
yang diberikan oleh $\omega(e_\lambda)=h_\lambda$ untuk setiap partisi $\lambda$.
Berdasarkan Proposisi~\ref{prop:fund_rel}, kita juga mempunyai
$\omega(h_\lambda)=e_\lambda$.
Dengan sedikit usaha tambahan, dapat ditunjukkan bahwa
$\omega(p_\lambda)=\epsilon_\lambda p_\lambda$
dan
$\omega(s_\lambda)=s_{\lambda^\dag}$.
Pengamatan terakhir ini menunjukkan bahwa $\omega$ adalah isometri
terhadap hasil kali dalam Hall:
$( \omega(f), \omega(g) ) = (f, g)$ untuk

semua $f,g \in \mathsf{SF}$.

\section{Representasi Grup Simetrik}
Misalkan $S_n$ grup simetrik pada $n$ huruf.
Ingat bahwa kelas-kelas konjugasi dari $S_n$ berkorespondensi bijektif dengan
partisi dari $n$ melalui tipe siklusnya.
Untuk partisi $\lambda \vdash n$ dan fungsi kelas
$f \in C(S_n)$, kita mendefinisikan $f(\lambda):=f(\sigma)$, dengan
$\sigma \in S_n$ sembarang permutasi bertipe siklus $\lambda$.
Misalkan $\chi_\lambda \in C(S_n)$ menyatakan fungsi karakteristik
\[
\chi_\lambda(\sigma) :=
\begin{cases}
1 & \textrm{jika $\sigma$ mempunyai tipe siklus $\lambda$,}\\
0 & \textrm{selainnya.}\\
\end{cases}
\]

\begin{theorem}
Terdapat isomorfisme grup isometrik
\[
\operatorname{ch} : R(S_n) \to \mathsf{SF}^n,
\]
yang disebut \emph{pemetaan karakteristik}, dan didefinisikan oleh
\[
\operatorname{ch}(f) :=
\sum_{\lambda \vdash n} \frac{f(\lambda)}{z_\lambda} p_{\lambda}.
\]
Jika dipandang sebagai fungsi kelas, nilai $f \in R(S_n)$
dapat dihitung melalui
\[
f(\lambda) = ( \operatorname{ch}(f), p_\lambda ).
\]
Fungsi karakteristik $\chi_\lambda \in C(S_n)$
bersesuaian dengan $\frac{p_{\lambda}}{z_\lambda}$ dalam
$\mathsf{SF}^n_{\mathbb{Q}}$.
Karakter tak tereduksi dalam $R(S_n)$
tepat merupakan $\chi^\lambda := \operatorname{chi}(s_\lambda)$
untuk $\lambda \vdash n$.
Karakter trivial $\chi^{(n)}$ bersesuaian dengan $h_n=s_n$, dan
karakter tanda $\chi^{(1^n)}$ bersesuaian dengan $e_n=s_{(1)^n}$.
Involusi $\omega : \mathsf{SF}^n \to \mathsf{SF}^n$
bersesuaian dengan operasi tensor oleh karakter tanda.

\end{theorem}
\begin{proof}
Kita hanya menunjukkan beberapa fakta menarik yang mudah dilihat.
Pertama, kita mendefinisikan \emph{pemetaan karakteristik}
$\operatorname{ch} : C(S_n) \to
\mathsf{SF}^n \otimes_{\mathbb{Z}} \mathbb{C}$,
karena belum jelas bahwa citra $R(S_n)$ bahkan
termuat dalam $\mathsf{SF}^n$.
Dengan mengingat $(p_\lambda,p_\mu)=z_\lambda \delta_{\mu \lambda}$, rumus
\[
f(\lambda) = ( \operatorname{ch}(f), p_\lambda )
\]
langsung diperoleh. Ini menunjukkan bahwa $\operatorname{ch}$ bijektif
(pada ruang vektor kompleks).

Korespondensi antara
$\chi_\lambda$ dan $\frac{p_{\lambda}}{z_\lambda}$
kini mengikuti dari pengamatan bahwa
\[
\chi_\lambda(\mu) = \delta_{\lambda \mu} =
\left(\frac{p_{\lambda}}{z_\lambda}, p_\mu\right)
\]
untuk semua $\lambda$ dan $\mu$.

Sekarang, ingat bahwa banyaknya elemen dalam $S_n$ bertipe siklus
$\lambda$ diberikan oleh $\frac{n!}{z_\lambda}$. Jadi
\[
(\chi_\lambda, \chi_\mu) =
\begin{cases}
\frac{1}{z_\lambda} & \textrm{jika $\lambda=\mu$,}\\
0 & \textrm{selainnya.}\\
\end{cases}
\]
Namun, ini tepat sama dengan $(p_\lambda,p_\mu)$,
sehingga $\operatorname{ch}$ adalah isometri.

Kita tidak menyertakan bukti pernyataan-pernyataan sisanya
(lihat \cite[\S{I.7}]{Macdonald} atau \cite[\S{7.18}]{Stanley2}).
\end{proof}

Korolari langsung dari teorema di atas adalah bahwa
tabel karakter $S_n$ tepat sama dengan
matriks perubahan basis antara $\{ p_\lambda \}$
dan $\{ s_\lambda \}$ dalam $\mathsf{SF}_n$.
Jika entri-entri tabel karakter $S_n$ dinotasikan
$\chi_\mu^\lambda$, maka
\[
\chi_\lambda = \sum_{\mu} \chi_\lambda^\mu \chi^\mu
\]
ekuivalen dengan
\[
p_\lambda = \sum_{\mu} \chi_\lambda^\mu s_\mu .
\]
Kini dapat dilihat bahwa entri-entri matriks pada Contoh~\ref{ex:transSP}
tepat merupakan entri-entri tabel karakter untuk $S_2$, $S_3$, dan
$S_4$.

\begin{definition}
Representasi tak tereduksi $V_\lambda$ dari $S_n$ yang bersesuaian dengan
$\chi^\lambda$ adalah \emph{modul Specht} yang terkait dengan $\lambda$.
\end{definition}

\begin{proposition}
$\displaystyle \dim_{\mathbb{C}} V_\lambda = f^\lambda
=\frac{n!}{\prod h(i,j)}$
\end{proposition}

\begin{proof}
Jika $f \in R(S_n)$ adalah karakter dari $V_\lambda$,
maka evaluasi pada identitas memberikan
\[
f(e)=f(1^d)= ( \operatorname{ch}(f), p_{(1^d)} )
= (s_\lambda, h_{(1^d)}) = K_{\lambda(1)^f}=f^\lambda.
\]
Sekarang gunakan Teorema~\ref{thm:hook_length}.
\end{proof}

Dengan mendefinisikan grup bergradasi
\[
R(S_\ast) := \bigoplus_{n \ge 0} R(S_n)
\]
kita memperoleh pemetaan karakteristik
\[
\operatorname{ch} : R(S_\ast) \to \mathsf{SF}
\]
dengan menjumlahkan komponen-komponen bergradasi.
Namun, perkalian yang sudah ada pada setiap $R(S_n)$
\emph{tidak} bersesuaian dengan perkalian pada $\mathsf{SF}$.
Kita mendefinisikan perkalian pada $R(S_\ast)$ yang menjadikan
$\operatorname{ch}$ suatu isomorfisme gelanggang bergradasi.

Terdapat pembenaman alami $S_n \times S_m \hookrightarrow S_{n+m}$;
ada banyak cara melakukannya, tetapi semuanya saling konjugat.
Untuk representasi $\rho$ dari $S_n$, kita memperoleh representasi
$\widetilde{\rho}$ dari $S_n \times S_m$ dari $\rho$ melalui prakomposisi
dengan proyeksi pertama. Serupa dengan itu, dari representasi $\sigma$
dari $S_m$, kita memperoleh representasi $\widetilde{\sigma}$ dari
$S_n \times S_m$.

Misalkan $R(S_n)$ adalah gelanggang representasi grup simetrik $S_n$. 
Kita mendefinisikan perkalian bilinear
\[ \boxtimes : R(S_n) \times R(S_m) \to R(S_{n+m}) \]
melalui
\[
\rho \boxtimes \sigma := \operatorname{Ind}_{S_n \times S_m}^{S_{n+m}}
(\widetilde{\rho} \otimes \widetilde{\sigma}).
\]
Ini menjadikan $R(S_\ast)$ gelanggang komutatif bergradasi,
dan $\operatorname{ch}$ merupakan isomorfisme gelanggang bergradasi.

\subsection{Deskripsi Eksplisit Modul Specht}

Modul Specht juga dapat dideskripsikan tanpa
mengacu pada fungsi simetris
(lihat \cite[\S{4}]{FultonHarris} dan \cite[\S{5.11-17}]{Etingof}).

Di sini kita benar-benar mengonstruksi setiap modul Specht $V_\lambda$
sebagai subrepresentasi dari representasi reguler $S_n$.
Perhatikan bahwa komponen isotipik yang terkait dengan $V_\lambda$ biasanya mempunyai
multiplisitas lebih besar dari $1$, sehingga wajar jika konstruksinya
bergantung pada pilihan sembarang.

\begin{definition}
Suatu \emph{tableau Young} $T$ adalah diagram Young untuk partisi $\lambda$
dari $n$, dengan setiap bilangan dalam $\{1,\ldots,n\}$ ditempatkan pada tepat satu
kotak. (Perhatikan bahwa ini bukan kasus khusus maupun perumuman
dari tableau Young standar dan semistandar yang dibahas di atas.)
Untuk tableau Young $T$, kita mendefinisikan $P_T$ sebagai
subgrup dari $S_n$ yang hanya mempermutasikan bilangan-bilangan dalam setiap baris $T$
dan $Q_T$ sebagai subgrup yang hanya mempermutasikan bilangan-bilangan dalam setiap
kolom $T$.  
\end{definition}

\begin{example}
Berikut adalah tableau Young untuk partisi $322$:
\begin{center}
\begin{ytableau}
3 & 1 & 5 \\
6 & 2 \\
4 & 7
\end{ytableau}
\end{center}
Kita mempunyai
\[
P_T = S_{\{1,3,5\}} \times S_{\{2,6\}} \times S_{\{4,7\}}
\textrm{ dan }
Q_T = S_{\{3,4,6\}} \times S_{\{1,2,7\}}
\]
dengan $S_X$ grup permutasi dari himpunan $X$. 
\end{example}

Jika $\lambda = (\lambda_1,\ldots,\lambda_r)$ suatu partisi,
definisikan $S_\lambda := S_{\lambda_1} \times \cdots \times S_{\lambda_r}$.
Perhatikan bahwa $P_T \cong S_\lambda$
dan $Q_T \cong S_{\lambda^\dag}$.
Subgrup $P_T$ dan $Q_T$ bergantung pada pilihan tableau $T$, tetapi
kelas konjugasinya hanya bergantung pada partisi yang bersesuaian,
$\lambda$.

Untuk tableau Young $T$ dan representasi $W$ dari $S_n$,
definisikan \emph{proyektor Young}:
\[
a_T := \frac{1}{|P_T|}\sum_{\sigma \in P_T} \sigma
\textrm{ dan }
b_T := \frac{1}{|Q_T|}\sum_{\sigma \in Q_T} \epsilon(\sigma)\sigma
\]
yang mudah dilihat sebagai proyeksi dalam $\operatorname{End}(W)$.
\emph{Penyimetris Young} adalah endomorfisme $c_T = a_T \circ b_T$.
Dapat diperiksa bahwa $c_T = b_T \circ a_T$.

\begin{theorem}
Andaikan $\lambda$ partisi dari $n$.
Jika $V_{reg}$ adalah representasi reguler dan $T$ suatu tableau Young
untuk $\lambda$, maka citra $c_T(V_{reg})$ isomorfik dengan
modul Specht $V_\lambda$.
\end{theorem}

\begin{proof}
Kita hanya membuat sketsanya.
Pertama, misalkan $U_T = a_T(V_{reg})$ dan $W_T = b_T(V_{reg})$.
Kita mengamati bahwa
\[
U_T \cong \operatorname{Ind}_{P_T}^{S_n} 1_{P_T}
\]
dan
\[
W_T \cong \operatorname{Ind}_{Q_T}^{S_n} \epsilon_{Q_T}
\]
dengan $1_{P_T}$ representasi trivial dan $\epsilon_{Q_T}$ representasi
tanda.
Dengan kata lain, karakteristik $U_T$ adalah
\[
\operatorname{ch}(1 \boxtimes \cdots \boxtimes 1)
= h_{\lambda_1} \cdots h_{\lambda_r} = h_\lambda,
\]
sedangkan karakteristik $W_T$ adalah
\[
\operatorname{ch}(\epsilon \boxtimes \cdots \boxtimes \epsilon)
= e_{\lambda^\dag_1} \cdots e_{\lambda^\dag_r} = e_{\lambda^\dag}.
\]
Dengan menuliskan $h_\lambda$ dan $e_{\lambda^\dag}$ dalam basis Schur,
kita melihat bahwa $s_\lambda$ adalah satu-satunya vektor basis dengan entri tak nol
pada keduanya.
Selain itu, $(s_\lambda,h_\lambda)=(s_\lambda,e_{\lambda^\dag})=1$.
Jadi citra $c_T = a_T \circ b_T$ isomorfik dengan modul Specht
$V_\lambda$, asalkan citranya tak nol.

Misalkan $\{ v_\sigma \}_{\sigma \in S_n}$ adalah basis standar bagi representasi
reguler $V_{reg}$.
Kita mempunyai deskripsi konkret
\[
c_T(v_{1})  = \frac{1}{|P_T| |Q_T|}\sum_{\sigma \in P_T}\sum_{\tau \in Q_T}
\epsilon(\tau) v_{\sigma\tau} .
\]
Karena $P_T \cap Q_T = \{1\}$, semua $v_{\sigma\tau}$ berbeda.
Jadi $c_T \ne 0$ sebagaimana diinginkan.
\end{proof}

\section{Representasi dari $\operatorname{GL}(V)$}

Funktor Schur dideskripsikan secara langsung dalam
\cite[\S{6}]{FultonHarris} dan \cite[\S{5.19,5.20-23}]{Etingof}.
Uraian yang sedikit lebih canggih dapat ditemukan dalam
\cite[Appendix 7.2]{Stanley2} dan \cite[Appendix I.A]{Macdonald}.

\begin{definition}
Untuk ruang vektor kompleks berdimensi $n$ yang dinotasikan $V$,
suatu \emph{representasi polinomial} (berturut-turut, \emph{representasi rasional}) dari
$\operatorname{GL}(V)$ adalah representasi
\[
\rho : \operatorname{GL}(V) \to \operatorname{GL}_N(\mathbb{C})
\]
dengan entri-entri $\rho(g)$ merupakan fungsi polinomial (berturut-turut, rasional)
dari entri-entri $g$.
\end{definition}

\begin{remark}
Bagi pembaca yang mengenal geometri aljabar,
representasi rasional merupakan pemetaan rasional pada gelanggang matriks
yang memuatnya, $\mathbb{C}^{n^2}$, tetapi bersifat reguler pada himpunan terbuka
$\operatorname{GL}_n(\mathbb{C})$.
Representasi polinomial bersifat reguler pada seluruh
$\mathbb{C}^{n^2}$.
Perlu berhati-hati: sebagian penulis memakai ``representasi polinomial'' sebagai sinonim
representasi rasional.
\end{remark}

\begin{example}
Jika $V = \mathbb{C}^2$, maka $\mathcal{S}^2(V) \cong \mathbb{C}^3$
mempunyai aksi alami $\operatorname{GL}(V)$.
Memang, dengan memilih basis $\{x,y\}$ bagi $V$ dan $\{x^2,xy,y^2\}$
bagi $\mathcal{S}^2(V)$, kita memperoleh representasi
$\rho : \operatorname{GL}_2(k) \to \operatorname{GL}_3(k)$
dengan
\[
\rho
\begin{pmatrix} a & b \\ c & d \end{pmatrix}
=
\begin{pmatrix}
a^2 & 2ab & b^2\\
ac & ad+bc & bd\\
c^2 & 2cd & d^2
\end{pmatrix} 
\]
sebagai deskripsi eksplisit.
\end{example}

\begin{example}
Jika $V$ berdimensi $n$, ingat bahwa $\Lambda^n(V)$ berdimensi $1$.
Representasi alami
$\rho: \operatorname{GL}(V) \to \operatorname{GL}(\Lambda^n(V)) \cong k$
adalah determinan.
\end{example}

Secara lebih umum, $\mathcal{S}^k(V)$ dan $\Lambda^k(V)$
mempunyai aksi kanonik $\operatorname{GL}(V)$ untuk semua $k \ge 0$.

\begin{example}
Tinjau representasi
$\rho : \operatorname{GL}_1(\mathbb{C}) \to
\operatorname{GL}_1(\mathbb{C})$
yang diberikan oleh
\[
\rho(z) = \frac{z}{|z|} .
\]
Representasi ini bukan representasi polinomial maupun rasional.
\end{example}

\begin{example}
\emph{Representasi dual} $V^\vee$ dari $V$ adalah aksi alami
$\operatorname{GL}(V)$ pada $V^\vee$.
Dengan memilih suatu basis beserta basis dualnya,
representasi dual bersesuaian dengan mengambil invers transpos
$(A^T)^{-1}$
dari matriks $A$.
Ini bukan representasi polinomial karena entri-entri $(A^T)^{-1}$ mempunyai
penyebut.
Namun, ingat bahwa matriks invers $A^{-1}$ berbentuk
$\det(A)\operatorname{adj}(A)$ dengan
$\operatorname{adj}(A)$ adalah \emph{matriks adjoin}.
Entri-entri matriks adjoin merupakan fungsi polinomial dari
entri-entri matriks semula.
Karena $(A^T)^{-1}=\operatorname{adj}(A)^T \det(A)^{-1}$,
kita menyimpulkan bahwa $V^\vee$ adalah representasi rasional.
\end{example}

Untungnya, perbedaan antara representasi polinomial dan rasional
mudah dikarakterisasi. Memang, penyebut yang muncul
dalam fungsi rasional hanya dapat berupa pangkat determinan; jika tidak,
representasinya tidak terdefinisi di setiap titik
$\operatorname{GL}(V)$.

\begin{proposition}
Jika $\rho$ representasi rasional dari $\operatorname{GL}(V)$,
maka terdapat bilangan bulat tak negatif minimal tunggal $m$
dan representasi polinomial tunggal $\sigma$
sedemikian sehingga
berlaku
\[
\rho(g) = \frac{\sigma(g)}{\det(g)^m}
\]
untuk semua $g \in \operatorname{GL}(V)$. 
\end{proposition}

Untuk ruang vektor kompleks berdimensi $n$ yang dinotasikan $V$,
ingat bahwa terdapat aksi alami $S_d$ pada pangkat tensor
$V^{\otimes d}$ yang memenuhi
\[
\sigma(v_1 \otimes \cdots \otimes v_d) =
\sigma(v_{\sigma^{-1}(1)} \otimes \cdots \otimes v_{\sigma^{-1}(d)})
\]
untuk $\sigma \in S_d$.
Terdapat pula aksi ``diagonal'' alami dari $\operatorname{GL}_n(V)$
pada $V^{\otimes d}$ yang memenuhi
\[
g(v_1 \otimes \cdots \otimes v_d) = g(v_1) \otimes \cdots \otimes g(v_d)
\]
untuk semua $g \in \operatorname{GL}(V)$.
Kedua aksi ini jelas komutatif, sehingga kita memperoleh aksi
$S_d \times \operatorname{GL}(V)$ pada $V^{\otimes d}$.

\begin{definition}
Untuk partisi $\lambda \vdash d$, kita
mempunyai \emph{funktor Schur} $\mathbb{S}_\lambda$,
yang memetakan ruang vektor $W$ ke
representasi-$\operatorname{GL}(W)$ berikut:
\[
\mathbb{S}_\lambda(W) :=
\operatorname{Hom}^{S_d}_{\mathbb{C}}(V_\lambda,W^{\otimes d}) \ .
\]
\end{definition}

\begin{remark}
Bagi pembaca yang mengenal teori kategori,
setiap funktor Schur $\mathbb{S}_\lambda$ adalah funktor kovarian
dari kategori ruang vektor berdimensi hingga ke dirinya sendiri.
Ini menyiratkan bahwa jika $f : W \to U$ suatu pemetaan linear,
maka kita juga mempunyai pemetaan linear
\[
\mathbb{S}_\lambda(f) : \mathbb{S}_\lambda(W) \to \mathbb{S}_\lambda(U).
\]
Fakta bahwa objek-objek ini merupakan representasi $\operatorname{GL}(W)$ dapat
dipandang sebagai akibat langsung dari sifatnya sebagai funktor.
(Khususnya, funktor Schur dapat didefinisikan untuk
kategori tensor yang lebih umum.)
\end{remark}

\begin{theorem}[Schur--Weyl]
Misalkan $W$ ruang vektor kompleks berdimensi $n$
dan misalkan $d \ge 0$.
Sebagai representasi-$S_d \times \operatorname{GL}(W)$, terdapat
dekomposisi kanonik
\[
W^{\otimes d} =
\bigoplus_{\substack{\lambda \vdash d\\ \ell(\lambda) \le n}}
V_\lambda \otimes \mathbb{S}_\lambda(W)
\]
dengan $V_\lambda$ modul Specht dari $\lambda$, dan setiap
\[
\mathbb{S}_\lambda(W) :=
\operatorname{Hom}^{S_d}_{\mathbb{C}}(V_\lambda,W^{\otimes d})
\]
merupakan representasi tak tereduksi dari $\operatorname{GL}(W)$.
\end{theorem}

\begin{theorem}
Kelas isomorfisme semua representasi \emph{polinomial}
tak tereduksi
dari $\operatorname{GL}(W)$ tepat merupakan $\mathbb{S}_\lambda(W)$
dengan $\ell(\lambda) \le d$.
Selain itu, jika $g \in \operatorname{GL}(W)$ mempunyai nilai-nilai eigen
$\alpha_1,\ldots, \alpha_n$, maka
\[
\operatorname{tr}(\mathbb{S}_\lambda(g)) = s_\lambda(\alpha_1,\alpha_2,\cdots,
\alpha_n,0,0,\cdots)
\]
sehingga $s_\lambda$ adalah karakter dari $\mathbb{S}_\lambda$.
\end{theorem}

Dengan mengevaluasi $\mathbb{S}_\lambda$ pada identitas, kita dapat menentukan
dimensi representasi tak tereduksi yang diperoleh dari funktor
Schur.

\begin{corollary}[Rumus Dimensi Weyl]
Jika $W$ ruang vektor berdimensi $n$ dan $\lambda$ suatu partisi.
Jika $\ell(\lambda) \le n$, maka
\[
\dim\ \mathbb{S}_\lambda(W) = s_\lambda(1,1,\ldots,1,0,\ldots)
= \prod_{1 \le i < j \le n} \frac{\lambda_i-\lambda_j +  j - i}{j-i}.
\]
Jika $\ell(\lambda) > n$, maka $\dim\ \mathbb{S}_\lambda(W) = 0$.
\end{corollary}

\begin{example}
Karena $V_{(d)}$ adalah representasi trivial dari $S_d$,
kita memperoleh
\[
\mathbb{S}_{(d)}(W) = \operatorname{Hom}^{S_d}_{\mathbb{C}}(\mathbb{C},W^{\otimes d})
\cong \operatorname{Sym}^d(W) \cong \mathcal{S}^d(W) .
\]
\end{example}

\begin{example}
Karena $V_{(1^d)}$ adalah representasi tanda dari $S_d$, yakni $\epsilon$,
kita memperoleh
\[
\mathbb{S}_{(1^d)}(W)L = \operatorname{Hom}^{S_d}_{\mathbb{C}}(\epsilon,W^{\otimes d})
\cong \operatorname{Alt}^d(W) \cong \Lambda^d(W) .
\]
Perhatikan bahwa $\mathbb{S}_{(1^d)}(W)=0$ ketika $\dim(W) > d$.
\end{example}

\bibliographystyle{alpha}
\bibliography{rep_theory}

\end{document}


